% GENERATED FILE. Edit canonical lessons, then run npm run generate:books.
% canonical-source-hash: fnv1a64:c18ae14e5427c680
% canonical-lessons: LA-C265-sicut, LA-C265-velut, LA-C265-utinam, LA-C265-immo, LA-C265-saltem

\chapter{Even as, As, If only, On the contrary, At least}
\label{ch:la-even-as-as-if-only-on-the-contrary-at-least}
\hlchaptermodality{265}

\begin{chapteropening}
\textbf{By the end of this chapter:} \emph{I can say even as, as, if only, on the contrary and at least.}

Complete the last lesson of chapter 265: I can say even as, as, if only, on the contrary and at least.
\end{chapteropening}

\section[sīcut]{sīcut — even as, in the same way as}

\label{lesson:LA-C265-sicut}



\begin{quote}
Before the new one: say the Latin for upwards, then the Latin for downwards.
\end{quote}

\subsection*{You'll want to know: sīcut}
\textbf{sīcut} — \textquotedblleft{}even as, in the same way as\textquotedblright{}.

\begin{cousinweb}
\textbf{sīcut} is \emph{sīc}, \textquotedblleft{}so\textquotedblright{}, with \emph{ut}, \textquotedblleft{}as\textquotedblright{} (Lewis \& Short): even as.
\end{cousinweb}

\begin{grammarlens}[title={What you've built}]
One more word for this chapter.
\end{grammarlens}

\subsection*{Your turn}
\begin{itemize}
\raggedright
  \item \emph{Say it:} \emph{sīcut}
  \item \emph{Say it:} \emph{sīcut}, once more
  \item \emph{Say it:} say \emph{sīcut}
  \item \emph{Recall:} say the Latin for upwards, then the Latin for downwards, then say \emph{sīcut} again
\end{itemize}

\begin{tcolorbox}[breakable,colback=teal!4,colframe=teal!35!black,title={Before you move on}]
What does sīcut mean? (\textquotedblleft{}Even as, in the same way as\textquotedblright{}.) Say it once more.
\end{tcolorbox}
\section[velut]{velut — as, like}

\label{lesson:LA-C265-velut}



\begin{quote}
Before the new one: say the Latin for downwards, then the Latin for even as.
\end{quote}

\subsection*{You'll want to know: velut}
\textbf{velut} — \textquotedblleft{}as, like\textquotedblright{}.

\begin{cousinweb}
\textbf{velut} is \emph{vel}, \textquotedblleft{}or, even\textquotedblright{}, with \emph{ut}, \textquotedblleft{}as\textquotedblright{} (Lewis \& Short): even as, like.
\end{cousinweb}

\begin{grammarlens}[title={What you've built}]
One more word for this chapter.
\end{grammarlens}

\subsection*{Your turn}
\begin{itemize}
\raggedright
  \item \emph{Say it:} \emph{velut}
  \item \emph{Say it:} \emph{velut}, once more
  \item \emph{Say it:} say \emph{velut}
  \item \emph{Recall:} say the Latin for downwards, then the Latin for even as, then say \emph{velut} again
\end{itemize}

\begin{tcolorbox}[breakable,colback=teal!4,colframe=teal!35!black,title={Before you move on}]
What does velut mean? (\textquotedblleft{}As, like\textquotedblright{}.) Say it once more.
\end{tcolorbox}
\section[utinam]{utinam — if only}

\label{lesson:LA-C265-utinam}



\begin{quote}
Before the new one: say the Latin for even as, then the Latin for as.
\end{quote}

\subsection*{You'll want to know: utinam}
\textbf{utinam} — \textquotedblleft{}if only\textquotedblright{}.

\begin{cousinweb}
\textbf{utinam} is \emph{ut}, \textquotedblleft{}how, that\textquotedblright{}, with the particle \emph{nam}, \textquotedblleft{}for, indeed\textquotedblright{} (EtymDB 2.1): oh that, if only.
\end{cousinweb}

\begin{grammarlens}[title={What you've built}]
One more word for this chapter.
\end{grammarlens}

\subsection*{Your turn}
\begin{itemize}
\raggedright
  \item \emph{Say it:} \emph{utinam}
  \item \emph{Say it:} \emph{utinam}, once more
  \item \emph{Say it:} say \emph{utinam}
  \item \emph{Recall:} say the Latin for even as, then the Latin for as, then say \emph{utinam} again
\end{itemize}

\begin{tcolorbox}[breakable,colback=teal!4,colframe=teal!35!black,title={Before you move on}]
What does utinam mean? (\textquotedblleft{}If only\textquotedblright{}.) Say it once more.
\end{tcolorbox}
\section[immō]{immō — on the contrary}

\label{lesson:LA-C265-immo}



\begin{quote}
Before the new one: say the Latin for as, then the Latin for if only.
\end{quote}

\subsection*{You'll want to know: immō}
\textbf{immō} — \textquotedblleft{}on the contrary\textquotedblright{}.

\begin{cousinweb}
\textbf{immō} has no sure origin; Lewis \& Short suggest a superlative form built on \emph{in}, \textquotedblleft{}in\textquotedblright{}, like \emph{summus}, \textquotedblleft{}highest\textquotedblright{}, but that is uncertain.
\end{cousinweb}

\begin{grammarlens}[title={What you've built}]
One more word for this chapter.
\end{grammarlens}

\subsection*{Your turn}
\begin{itemize}
\raggedright
  \item \emph{Say it:} \emph{immō}
  \item \emph{Say it:} \emph{immō}, once more
  \item \emph{Say it:} say \emph{immō}
  \item \emph{Recall:} say the Latin for as, then the Latin for if only, then say \emph{immō} again
\end{itemize}

\begin{tcolorbox}[breakable,colback=teal!4,colframe=teal!35!black,title={Before you move on}]
What does immō mean? (\textquotedblleft{}On the contrary\textquotedblright{}.) Say it once more.
\end{tcolorbox}
\section[saltem]{saltem — at least}

\label{lesson:LA-C265-saltem}



\begin{quote}
Before the new one: say the Latin for if only, then the Latin for on the contrary.
\end{quote}

\subsection*{You'll want to know: saltem}
\textbf{saltem} — \textquotedblleft{}at least\textquotedblright{}.

\begin{cousinweb}
\textbf{saltem} has no sure origin; Lewis \& Short tie it to \emph{salvus}, \textquotedblleft{}safe\textquotedblright{}, but that is uncertain.
\end{cousinweb}

\begin{grammarlens}[title={What you've built}]
One more word for this chapter.
\end{grammarlens}

\subsection*{Your turn}
\begin{itemize}
\raggedright
  \item \emph{Say it:} \emph{saltem}
  \item \emph{Say it:} \emph{saltem}, once more
  \item \emph{Say it:} say \emph{saltem}
  \item \emph{Recall:} say the Latin for if only, then the Latin for on the contrary, then say \emph{saltem} again
\end{itemize}

\begin{tcolorbox}[breakable,colback=teal!4,colframe=teal!35!black,title={Before you move on}]
What does saltem mean? (\textquotedblleft{}At least\textquotedblright{}.) Say it once more.
\end{tcolorbox}
